\documentclass{article}
\usepackage[utf8]{inputenc}
\usepackage{hyperref}

\let\originalfootnotesize\footnotesize
\let\oldverbatim\verbatim
\let\oldendverbatim\endverbatim
\renewenvironment{verbatim}{\originalfootnotesize\oldverbatim}{\oldendverbatim}

\title{Incident Response Management}
\author{ERP/1 SRE Team}
\date{\today}

\begin{document}

\maketitle

\begin{abstract}
This article presents a lightweight, Git-native methodology for managing incident response
playbooks within the ERP/1 Kubernetes platform. By maintaining all documentation
in version-controlled Markdown files, the approach achieves maximal simplicity, auditability,
traceability, and seamless integration with GitOps principles.
\end{abstract}

\tableofcontents

\newpage
\section{Introduction}

In Site Reliability Engineering, playbooks are critical for reducing Mean Time To Recovery (MTTR).
Google SRE practices demonstrate that well-documented playbooks can improve incident resolution times by
compared to ad-hoc responses. For the ERP/1 platform, we adopted a pure git strategy \footnote{\url{https://github.com/erpuno/cd/tree/main/playbooks}}:

\begin{itemize}
    \item No heavy tools or databases
    \item Full version history and collaboration via pull requests
    \item Seamless integration with ArgoCD and the existing CD pipeline
    \item Easy search and rendering on GitHub
\end{itemize}

\subsection{Core Principles from Google SRE}

\begin{itemize}
    \item \textbf{Playbooks dramatically improve MTTR} — Google observed roughly 3x faster resolution when on-call engineers follow documented playbooks instead of “winging it.”
    \item \textbf{Preparation beats heroics} — Invest in up-to-date, actionable documentation + regular practice (e.g., “Wheel of Misfortune” drills).
    \item \textbf{Blameless culture} — Focus postmortems on systems and processes, not people.
    \item \textbf{Automate where possible} — But keep humans in the loop for judgment.
\end{itemize}

\subsection{The 5 A’s Framework for Excellent Playbooks}

A widely recommended structure for high-quality playbooks includes:

\begin{enumerate}
    \item \textbf{Actionable} — Every step is a clear command/checklist (not prose).
    \item \textbf{Accessible} — Available where responders work (Git, linked from alerts).
    \item \textbf{Accurate} — Regularly reviewed and tested.
    \item \textbf{Authoritative} — Single source of truth with ownership last update date.
    \item \textbf{Adaptable} — Easy to update; includes decision trees for variations.
\end{enumerate}

\newpage
\subsection{Key Structural Best Practices}

\begin{itemize}
    \item Standard Template per playbook — Title + Severity + Impact, Detection (alerts / SLIs),
          Immediate Actions, Diagnosis / Troubleshooting flow (with commands), Resolution steps, Escalation path, Prevention / Follow-up, Metadata.
    \item Keep it concise — Aim for scannable under stress. Use numbered steps, code blocks, and decision trees.
    \item Make it executable — Tools like Runme turn Markdown code blocks into one-click commands.
    \item Version control — Store in Git (Markdown) alongside code. Link to service diagrams, dashboards, and dependencies.
    \item Service-specific and Generic — Combine broad (e.g., “Resource Exhaustion”) and component-specific playbooks.
\end{itemize}

\subsection{Operational Best Practices}

\begin{itemize}
    \item Regular reviews \& testing — Quarterly or after every Sev1/2 incident. Test via chaos engineering and tabletop exercises.
    \item Integrate with tools — Link playbooks directly from Prometheus/Grafana alerts.
    \item Train \& practice — On-call rotations should include playbook familiarization and simulated incidents.
    \item Measure \& improve — Track playbook usage versus MTTR. Update based on real postmortems.
    \item Minimize toil — Automate repetitive steps where appropriate.
    \item Include architecture context — Diagrams, dependencies, recent changes, known issues.
\end{itemize}

\subsection{Common Pitfalls to Avoid}

\begin{itemize}
    \item Outdated steps (leads to worse outcomes).
    \item Too much narrative instead of commands.
    \item Buried in wikis (hard to find during incidents).
    \item No ownership or review process.
    \item Ignoring non-technical aspects (communication templates, stakeholder updates).
\end{itemize}

\section{ERP/1 Playbooks}

\subsection{Detailed Playbook Template Specification}

To satisfy all requirements of the 5 A’s framework and Google SRE best practices,
every playbook in the repository follows this standardized template:

\begin{itemize}
    \item \textbf{Title \& Metadata} — Clear title, severity level, impact statement, owner, last reviewed date.
    \item \textbf{Impact \& Scope} — Description of affected services/namespaces and potential business impact.
    \item \textbf{Detection} — Specific alerts, SLIs, or symptoms that trigger the playbook.
    \item \textbf{Immediate Actions} (First 5 Minutes) — High-priority steps to stabilize the situation.
    \item \textbf{Diagnosis} — Step-by-step troubleshooting with commands, logs, and decision trees.
    \item \textbf{Resolution} — Detailed remediation procedures, including rollback and verification steps.
    \item \textbf{Escalation Path} — Clear criteria and contacts for further escalation.
    \item \textbf{Prevention \& Follow-up} — Long-term improvements, chaos tests, and postmortem requirements.
    \item \textbf{References} — Links to architecture diagrams, dashboards, related playbooks, and playbooks.
\end{itemize}

This template ensures playbooks are concise yet comprehensive, scannable under pressure,
and fully executable (using code blocks). It balances generic patterns with service-specific guidance.

\subsection{Architecture}

The playbooks live in a dedicated \texttt{playbooks/} directory:

\begin{verbatim}
  README.md
  01-sev1-cluster-outage.md
  02-sev2-application-degradation.md
  03-sev2-resource-exhaustion.md
  04-sev2-gitops-network-stateful.md
\end{verbatim}

Each file follows a template optimized for on-call use under pressure.

\subsection{Best Practices Implemented}

\begin{enumerate}
    \item \textbf{Actionable} — Every playbook contains ready-to-run \texttt{kubectl}, \texttt{argocd}.
    \item \textbf{Versioned} — Git tracks every change with clear commit messages.
    \item \textbf{Searchable} — GitHub search + static site.
    \item \textbf{Executable} — Compatible for one-click execution.
    \item \textbf{Reviewed} — Changes via PRs with SRE approval.
\end{enumerate}

\subsection{Usage During Incidents}

\begin{enumerate}
    \item Open the relevant playbook directly from GitHub or local clone.
    \item Follow numbered steps.
    \item Document actions in the incident record.
    \item After resolution, update the playbook if new lessons were learned.
\end{enumerate}

\subsection{Advantages of Pure Git Approach}

\begin{itemize}
    \item Zero operational overhead
    \item Works offline
    \item Perfect audit trail for compliance (NIST)
    \item Natural fit for GitOps culture
    \item Easy to extend with automation (e.g., linking from Prometheus alerts)
\end{itemize}

\subsection{Future Enhancements}

\begin{itemize}
    \item Runme integration for executable blocks
    \item Automated testing of playbook steps
    \item AI-assisted playbook generation
\end{itemize}

\section{Conclusion}

Managing playbooks in pure Git represents the SRE ideal: simplicity, reliability,
and developer-friendly workflows. This approach keeps our focus on reliability rather than tool administration.

\end{document}
